% from ega1-9-fr.tex
\oldpage[I]{179}
\begin{corollary}[9.5.11]
\label{I.9.5.11-fr}
Sous les hypothèses de (\hyperref[I.9.5.10-fr]{9.5.10}), toute section de
$\mathscr{O}_{\overline{Y}}$ au-dessus d'un ouvert $V$ de $\overline{Y}$
qui est nulle dans $V\cap Y$ est nulle.
\end{corollary}

\begin{proof}
En vertu de (\hyperref[I.9.5.8-fr]{9.5.8}), on peut se ramener au cas où
$V=\overline{Y}$. Si l'on tient compte de ce que les sections de
$\mathscr{O}_{\overline{Y}}$ au-dessus de $\overline{Y}$ correspondent
canoniquement aux $\overline{Y}$-sections de
$\overline{Y}\otimes_{\mathbb{Z}}\mathbb{Z}[T]$
(\hyperref[I.3.3.15-fr]{3.3.15}) et de ce que ce dernier est séparé sur
$\overline{Y}$, le corollaire apparaît comme un cas particulier de
(\hyperref[I.9.5.6-fr]{9.5.6}).
\end{proof}

Lorsqu'il existe un plus petit sous-préschéma fermé $Y$ de $X$
majorant un sous-préschéma $Y$ de $X$, on dit que $Y'$ est
l'\emph{adhérence} de $Y$ dans $X$, lorsqu'il n'en résulte pas de
confusion.

\subsection{Faisceaux quasi-cohérents d'algèbres~; changement du faisceau
structural.}
\label{subsection:I.9.6-fr}

\begin{proposition}[9.6.1]
\label{I.9.6.1-fr}
Soient $X$ un préschéma, $\mathscr{B}$ une $\mathscr{O}_X$-Algèbre
quasi-cohérente (\hyperref[0.5.1.3-fr]{0, 5.1.3}). Pour qu'un
$\mathscr{B}$-Module $\mathscr{F}$ soit quasi-cohérent (sur l'espace
annelé $(X,\mathscr{B})$), il faut et il suffit que $\mathscr{F}$ soit un
$\mathscr{O}_X$-Module quasi-cohérent.
\end{proposition}

\begin{proof}
Comme la question est locale, on peut supposer $X$ affine d'anneau $A$, et
alors $\mathscr{B}=\widetilde{B}$, où $B$ est une $A$-algèbre
(\hyperref[I.1.4.3-fr]{1.4.3}). Si $\mathscr{F}$ est quasi-cohérent sur
l'espace annelé $(X,\mathscr{B})$, on peut aussi supposer que
$\mathscr{F}$ est le conoyau d'un $\mathscr{B}$-homomorphisme
$\mathscr{B}^{(I)}\to\mathscr{B}^{(J)}$~; comme cet homomorphisme est
aussi un $\mathscr{O}_X$-homomorphisme de $\mathscr{O}_X$-Modules, et que
$\mathscr{B}^{(I)}$ et $\mathscr{B}^{(J)}$ sont des
$\mathscr{O}_X$-Modules quasi-cohérents
(\hyperref[I.1.3.9-fr]{1.3.9, (ii)}), $\mathscr{F}$ est aussi un
$\mathscr{O}_X$-Module quasi-cohérent
(\hyperref[I.1.3.9-fr]{1.3.9, (i)}).

Inversement, si $\mathscr{F}$ est un $\mathscr{O}_X$-Module
quasi-cohérent, on a $\mathscr{F}=\widetilde{M}$, où $M$ est un $B$-module
(\hyperref[I.1.4.3-fr]{1.4.3})~; $M$ est isomorphe à un conoyau de
$B$-homomorphisme $B^{(I)}\to B^{(J)}$, donc $\mathscr{F}$ est un
$\mathscr{B}$-Module isomorphe au conoyau de l'homomorphisme correspondant
$\mathscr{B}^{(I)}\to\mathscr{B}^{(J)}$
(\hyperref[I.1.3.13-fr]{1.3.13}), ce qui achève la démonstration.
\end{proof}

En particulier, si $\mathscr{F}$ et $\mathscr{G}$ sont deux
$\mathscr{B}$-Modules quasi-cohérents,
$\mathscr{F}\otimes_{\mathscr{B}}\mathscr{G}$ est un
$\mathscr{B}$-Module quasi-cohérent~; il en est de même de
$\mathscr{H}\!om_{\mathscr{B}}(\mathscr{F},\mathscr{G})$ lorsqu'on
suppose en outre que $\mathscr{F}$ admet une présentation finie
(\hyperref[I.1.3.13-fr]{1.3.13}).

\begin{env}[9.6.2]
\label{I.9.6.2-fr}
Étant donné un préschéma $X$, nous dirons qu'une
$\mathscr{O}_X$-Algèbre quasi-cohérente $\mathscr{B}$ est de
\emph{type fini} si pour tout $x\in X$, il existe un voisinage ouvert
\emph{affine} $U$ de $x$ tel que $\Gamma(U,\mathscr{B})=B$ soit une
algèbre de type fini sur $\Gamma(U,\mathscr{O}_X)=A$. On a alors
$\mathscr{B}|U=\widetilde{B}$ et pour tout $f\in A$, la
$(\mathscr{O}_X|D(f))$-Algèbre induite $\mathscr{B}|D(f)$ est de type
fini, car elle est isomorphe à $\widetilde{B_f}$, et
$B_f=B\otimes_A A_f$ est évidemment une algèbre de type fini sur $A_f$.
Comme les $D(f)$ forment une base de la topologie de $U$, on en conclut que
si $\mathscr{B}$ est une $\mathscr{O}_X$-Algèbre quasi-cohérente de type
fini, pour tout ouvert $V$ de $X$, $\mathscr{B}|V$ est une
$(\mathscr{O}_X|V)$-Algèbre quasi-cohérente de type fini.
\end{env}

\begin{proposition}[9.6.3]
\label{I.9.6.3-fr}
Soit $X$ un préschéma localement noethérien. Toute
$\mathscr{O}_X$-Algèbre quasi-cohérente de type fini $\mathscr{B}$ est
alors un faisceau cohérent d'anneaux
(\hyperref[0.5.3.7-fr]{0, 5.3.7}).
\end{proposition}

\begin{proof}
On peut de nouveau se limiter au cas où $X$ est un schéma affine d'anneau
noethérien $A$, et où $\mathscr{B}=\widetilde{B}$, $B$ étant une
$A$-algèbre de type fini~; $B$ est donc un anneau noethérien. Cela étant, il
faut prouver que le noyau $\mathscr{N}$ d'un
$\mathscr{B}$-homomorphisme $\mathscr{B}^m\to\mathscr{B}$ est un
$\mathscr{B}$-
